%
% Documentation for the Hotel Booking Manager project
% ===================================================
\documentclass[11pt,a4paper]{article}

\usepackage{thga-db}

\renewcommand{\thgaKurs}{Databases}
\renewcommand{\thgaDatum}{Summer Term 2026}

\begin{document}
\thispagestyle{empty}

\thgaTitel{Term Project -- Documentation}%
{Hotel Booking Manager}

\begin{infobox}
\textbf{Project:} Hotel Booking Manager \\
\textbf{Module:} Introduction to Database Management Systems -- THGA Bochum \\
\textbf{Stack:} PostgreSQL -- FastAPI -- Python/Tkinter \\
\textbf{Repository:} \url{https://github.com/abdelillah003/DBMS_10}
\end{infobox}

\tableofcontents
\vspace{1em}


% ============================================================
\section{Introduction and motivation}
% ============================================================

The Hotel Booking Manager is a database-backed desktop application for
managing room reservations in a small hotel. The system replaces
spreadsheets or handwritten reservation lists with a structured application
based on a relational PostgreSQL database.

The main users of the application are hotel reception staff. They can view
hotel rooms, inspect existing bookings, create new reservations, view booking
statistics, and create additional hotel services.

A relational database is suitable for this application because the data has
clear relationships. A booking must reference an existing guest and an
existing room. Rooms belong to room categories, and additional services can
be associated with bookings.

Primary keys, foreign keys, uniqueness constraints, validation rules, and
database-level constraints help to prevent inconsistent data.

The project consists of three main components:

\begin{itemize}
    \item a PostgreSQL database,
    \item a REST API implemented with FastAPI,
    \item a desktop frontend implemented with Python and Tkinter.
\end{itemize}


% ============================================================
\section{Data model}
% ============================================================

The database contains six main tables:

\begin{itemize}
    \item \texttt{guest}: stores information about hotel guests,
    \item \texttt{room\_category}: stores room categories such as Single,
          Double, and Suite,
    \item \texttt{room}: stores the individual hotel rooms,
    \item \texttt{booking}: stores room reservations,
    \item \texttt{service}: stores additional hotel services,
    \item \texttt{booking\_service}: connects bookings and services.
\end{itemize}

The main relationships are:

\begin{itemize}
    \item one room category can contain many rooms,
    \item one guest can have many bookings,
    \item one room can have many bookings over time,
    \item one booking can use multiple services,
    \item one service can be associated with multiple bookings.
\end{itemize}

The many-to-many relationship between bookings and services is resolved using
the \texttt{booking\_service} table.


\subsection{Relational schema}

The relational schema consists of the following tables:

\begin{verbatim}
guest
room_category
room
booking
service
booking_service
\end{verbatim}

Each table represents one entity or one relationship. Primary keys uniquely
identify records, while foreign keys represent the relationships between
tables.

For example, every booking references a guest and a room. Every room
references its room category. The \texttt{booking\_service} table references
both a booking and a service.

The schema is normalized to Third Normal Form (3NF). Each table stores
information about one subject, and non-key attributes depend on the key of
their table. Repeated information such as room category data is stored only
once and referenced through foreign keys. This reduces redundancy and helps
to keep the data consistent.


% ============================================================
\section{Database (PostgreSQL)}
% ============================================================

PostgreSQL is used as the relational database management system.

The database schema is stored in:

\begin{verbatim}
database/schema.sql
\end{verbatim}

Initial example data is stored in:

\begin{verbatim}
database/seed.sql
\end{verbatim}

The database uses primary keys and foreign keys to maintain relationships
between the tables. Additional constraints are used to reject invalid data.

An important business rule is that the same room must not be booked for
overlapping periods. This rule is implemented directly in PostgreSQL using
the exclusion constraint:

\begin{verbatim}
no_overlapping_bookings
\end{verbatim}

Because this protection is implemented at database level, conflicting
reservations are rejected even if they are submitted without using the
Tkinter frontend.

This is an important design decision because data integrity does not depend
only on application-side validation.


% ============================================================
\section{REST API (FastAPI)}
% ============================================================

The backend is implemented in Python using FastAPI. Communication with
PostgreSQL is handled using the \texttt{psycopg} driver.

The API contains the five main endpoints defined for the project:

\begin{center}
\begin{tabular}{lll}
\hline
\textbf{Method} & \textbf{Path} & \textbf{Purpose} \\
\hline
GET  & \texttt{/rooms} &
Room information \\
GET  & \texttt{/bookings} &
Booking information \\
GET  & \texttt{/statistics/bookings} &
Booking statistics \\
POST & \texttt{/bookings} &
Create a booking \\
POST & \texttt{/services} &
Create a service \\
\hline
\end{tabular}
\end{center}

The backend also provides the root endpoint:

\begin{verbatim}
GET /
\end{verbatim}

This endpoint can be used as a simple health check for the API.


\subsection{Read operations}

The read endpoints combine data from related tables.

The endpoint \texttt{GET /rooms} joins rooms with room categories. The
response includes information such as room ID, room number, floor, category,
capacity, and price per night.

The endpoint \texttt{GET /bookings} joins bookings with guests and rooms.
This allows the client to receive guest names and room numbers together with
the booking information.

The endpoint \texttt{GET /statistics/bookings} uses PostgreSQL aggregation
with \texttt{COUNT} and \texttt{FILTER}. It returns the total number of
bookings and the numbers of confirmed, checked-in, checked-out, and cancelled
bookings.


\subsection{Creating bookings}

The endpoint

\begin{verbatim}
POST /bookings
\end{verbatim}

creates a new booking.

The backend validates that the check-out date is after the check-in date. If
the dates are invalid, HTTP status \texttt{400 Bad Request} is returned.

If the specified guest or room does not exist, the request is also rejected
with HTTP status \texttt{400 Bad Request}.

If the PostgreSQL exclusion constraint detects an overlapping reservation for
the same room, the API returns:

\begin{verbatim}
409 Conflict
\end{verbatim}

A valid booking is created with HTTP status:

\begin{verbatim}
201 Created
\end{verbatim}


\subsection{Creating services}

The endpoint

\begin{verbatim}
POST /services
\end{verbatim}

creates a new hotel service.

Negative prices are rejected. Duplicate service names are also rejected. A
duplicate service name results in HTTP status \texttt{409 Conflict}.


\subsection{API authentication}

Write operations are protected with an X-API-Key.

The client sends the following HTTP header:

\begin{verbatim}
X-API-Key: hotel-booking-key
\end{verbatim}

The protected endpoints are:

\begin{verbatim}
POST /bookings
POST /services
\end{verbatim}

If the key is missing or incorrect, the API returns:

\begin{verbatim}
401 Unauthorized
\end{verbatim}

The GET endpoints remain accessible without authentication.

The authentication mechanism was tested using FastAPI's generated Swagger
interface. Requests without the key were rejected, while requests with the
correct key were accepted.


% ============================================================
\section{Frontend (Tkinter)}
% ============================================================

The desktop frontend is implemented in Python using Tkinter. It communicates
with the FastAPI backend through HTTP requests.

The application contains the following views:

\begin{itemize}
    \item Connection,
    \item Main Menu,
    \item Rooms,
    \item Bookings,
    \item Create Booking,
    \item Booking Statistics,
    \item Add Service.
\end{itemize}


\subsection{Connection}

The Connection view allows the user to enter the API URL and test whether the
backend is reachable.

For the local SSH-tunnel setup, the API URL is:

\begin{verbatim}
http://localhost:8003
\end{verbatim}


\subsection{Rooms}

The Rooms view retrieves its data from:

\begin{verbatim}
GET /rooms
\end{verbatim}

It displays the room ID, room number, floor, room category, capacity, and
price per night.

Displaying the database room ID is useful when a new booking is created,
because the booking form requires the room ID.


\subsection{Bookings}

The Bookings view uses:

\begin{verbatim}
GET /bookings
\end{verbatim}

It displays information including:

\begin{itemize}
    \item booking ID,
    \item guest name,
    \item room number,
    \item check-in date,
    \item check-out date,
    \item booking status.
\end{itemize}

The view also provides a refresh function so that newly created bookings can
be shown immediately.


\subsection{Create Booking}

The Create Booking view sends data to:

\begin{verbatim}
POST /bookings
\end{verbatim}

The user enters:

\begin{itemize}
    \item guest ID,
    \item room ID,
    \item check-in date,
    \item check-out date.
\end{itemize}

The frontend automatically includes the required X-API-Key in the HTTP
request.

Errors returned by the backend, including booking conflicts, are displayed
to the user.


\subsection{Booking Statistics}

The statistics view retrieves data from:

\begin{verbatim}
GET /statistics/bookings
\end{verbatim}

It displays:

\begin{itemize}
    \item total bookings,
    \item confirmed bookings,
    \item checked-in bookings,
    \item checked-out bookings,
    \item cancelled bookings.
\end{itemize}


\subsection{Add Service}

The Add Service view sends new service information to:

\begin{verbatim}
POST /services
\end{verbatim}

The frontend automatically includes the X-API-Key for this write request.


% ============================================================
\section{Operation and deployment}
% ============================================================

The PostgreSQL database and FastAPI backend run on the Linux server.

The backend can be started from the project directory with:

\begin{verbatim}
source .venv/bin/activate
python3 -m uvicorn backend.main:app --host 0.0.0.0 --port 8003
\end{verbatim}

FastAPI provides interactive Swagger documentation at:

\begin{verbatim}
http://localhost:8003/docs
\end{verbatim}

When the backend is running on the remote Linux server, an SSH tunnel can be
created from the Windows computer:

\begin{verbatim}
ssh -L 8003:localhost:8003 user@server
\end{verbatim}

The local port \texttt{8003} is then forwarded to port \texttt{8003} on the
server.

The Tkinter frontend runs on the Windows computer and connects to:

\begin{verbatim}
http://localhost:8003
\end{verbatim}

The frontend can be started with:

\begin{verbatim}
python main.py
\end{verbatim}

Inside the repository, the application components are stored in:

\begin{verbatim}
backend/main.py
frontend/main.py
database/schema.sql
database/seed.sql
\end{verbatim}


% ============================================================
\section{Testing}
% ============================================================

The system was tested at database, API, and frontend level.


\subsection{Database tests}

The database was tested with valid and invalid booking data.

An overlapping reservation for the same room was rejected by PostgreSQL.
This confirmed that the exclusion constraint works independently of the REST
API.


\subsection{API tests}

The FastAPI endpoints were tested using Swagger/OpenAPI.

The following behaviour was verified:

\begin{itemize}
    \item rooms can be retrieved successfully,
    \item bookings can be retrieved successfully,
    \item booking statistics can be retrieved successfully,
    \item valid bookings can be created,
    \item invalid date ranges are rejected,
    \item invalid guest or room references are rejected,
    \item overlapping bookings are rejected,
    \item hotel services can be created,
    \item duplicate service names are rejected,
    \item missing API keys are rejected,
    \item the correct API key allows write operations.
\end{itemize}


\subsection{Frontend tests}

The Tkinter frontend was tested against the running FastAPI backend.

A new booking was successfully created through the frontend and appeared in
the booking list afterward.

A new hotel service was also successfully created through the frontend.

These tests verify the complete communication path:

\begin{verbatim}
Tkinter frontend
       |
       v
FastAPI REST API
       |
       v
PostgreSQL database
\end{verbatim}


% ============================================================
\section{Reflection}
% ============================================================

The project demonstrates how a relational PostgreSQL database can be
integrated into a complete application consisting of a REST API and a desktop
user interface.

One of the most important design decisions was to enforce the rule against
overlapping room reservations directly in PostgreSQL. This provides stronger
data integrity than relying only on checks in the frontend.

FastAPI made it possible to implement a clear REST interface and provided
Swagger/OpenAPI documentation automatically. This was especially useful
during testing.

The separation between database, API, and frontend also makes the system
easier to understand and maintain. The Tkinter application does not access
PostgreSQL directly; all application requests go through the REST API.

The current implementation intentionally concentrates on the endpoints
defined in the project scope.

Possible future improvements include:

\begin{itemize}
    \item managing guests through the frontend,
    \item assigning existing services directly to bookings,
    \item changing booking statuses,
    \item searching specifically for available rooms,
    \item storing the API key in an environment variable,
    \item packaging the frontend for easier installation.
\end{itemize}

Overall, the Hotel Booking Manager fulfills its main objective: hotel
reservation data is stored in a normalized PostgreSQL database and accessed
through a FastAPI REST interface and a Tkinter desktop application.


\end{document}
